\documentclass{beamer}

\usepackage{amsmath, amssymb}
\usepackage{cleveref}

% -------------------------------------------------------
% Section header on every content slide (top-right corner)
% -------------------------------------------------------
\setbeamertemplate{headline}{%
  \begin{beamercolorbox}[wd=\paperwidth, ht=2.5ex, dp=1ex, rightskip=0.3cm]{section in head/foot}%
    \hfill\small\textit{\insertsectionhead}\hspace{0.3cm}
  \end{beamercolorbox}%
}

\title{Inference for Categorical Time Series with Missing Observations}
\author{Victor Minig}
\date{27.04.2026}

\begin{document}

\frame{\titlepage}

%------------------------------------------------
\begin{frame}{Outline}
\tableofcontents
\end{frame}

%================================================
\section{Categorical Time Series}
%================================================

\begin{frame}{Motivation}
\begin{itemize}
    \item Many real-world time series are \textbf{categorical}:
    \begin{itemize}
        \item Credit ratings (AAA, AA, ...)
        \item Weather states (sunny, rainy, ...)
        \item Survey responses (agree, neutral, disagree)
    \end{itemize}
    
    \vspace{0.3cm}
    
    \item Classical time series methods assume \textbf{numeric structure}
    
    \vspace{0.3cm}
    
    \item $\Rightarrow$ Need new tools for:
    \begin{itemize}
        \item Computing proper tests for marginal features 
        \item Measuring temporal dependence without bias 
    \end{itemize}
\end{itemize}
\end{frame}

%------------------------------------------------
\begin{frame}{Categorical Time Series}
\begin{block}{Definition}
A categorical time series $(X_t)$ takes values in a finite set
\[
S = \{s_0, \dots, s_m\}.
\]
\end{block}

\vspace{0.3cm}

\textbf{Two types of categorical data:}
\begin{itemize}
    \item \textbf{Nominal:} no natural ordering
    \item \textbf{Ordinal:} inherent ordering
\end{itemize}

\vspace{0.3cm}

\textbf{Key implication:}
\begin{itemize}
    \item Available statistical tools depend on structure of $S$
\end{itemize}
\end{frame}

%------------------------------------------------
\begin{frame}{Marginal Distributions}
\textbf{Central object: marginal probabilities}

\vspace{0.3cm}

\textbf{Nominal data:}
\[
p_i = P(X_t = s_i)
\]
\begin{itemize}
    \item Probability Mass Function (PMF)
\end{itemize}

\vspace{0.5cm}

\textbf{Ordinal data:}
\[
f_i = P(X_t \leq s_i)
\]
\begin{itemize}
    \item Cumulative Distribution Function (CDF)
\end{itemize}

\vspace{0.3cm}

\textbf{Key difference:}
\begin{itemize}
    \item Ordinal $\Rightarrow$ uses ordering
    \item Nominal $\Rightarrow$ no cumulative structure
\end{itemize}
\end{frame}

%------------------------------------------------
\begin{frame}{Location Measures}
\textbf{Nominal data:}
\begin{itemize}
    \item No ordering $\Rightarrow$ no mean/median
    \item Use \textbf{mode}:
\[
\text{Mod}(X_t) = \arg\max_{s_i} P(X_t = s_i)
\]
\end{itemize}

\vspace{0.5cm}

\textbf{Ordinal data:}
\begin{itemize}
    \item Ordering allows additional measures
    \item \textbf{Median:}
\[
\text{Med}(X_t) = \min \{ s_i : F(s_i) \geq 1/2 \}
\]
\end{itemize}

\vspace{0.3cm}

\textbf{Interpretation:}
\begin{itemize}
    \item Mode = most frequent category
    \item Median = central position in ordering
\end{itemize}
\end{frame}

%------------------------------------------------
\begin{frame}{Dispersion Measures}
\textbf{Different notions of variability}

\vspace{0.3cm}

Both indices are scaled to $[0,1]$, with 0 = minimal dispersion (mass on single category) and 1 = maximal dispersion (differs between types).

\vspace{0.4cm}

\textbf{Ordinal: Index of Ordinal Variation (IOV)}
\[
\text{IOV}(\boldsymbol{f}) = \frac{4}{m} \sum_{i=0}^{m-1} f_i (1 - f_i)
\]

\vspace{0.4cm}

\textbf{Nominal: Index of Qualitative Variation (IQV)}
\[
\text{IQV}(\boldsymbol{p}) = \frac{m+1}{m} \left(1 - \sum_{i=0}^{m} p_i^2 \right)
\]

\end{frame}

%------------------------------------------------
\begin{frame}{Stationarity}
\textbf{Why do we need stationarity?}

\begin{itemize}
    \item Stable probabilistic structure over time
    \item Enables estimation and asymptotics
\end{itemize}

\vspace{0.5cm}

\begin{block}{Strict Stationarity}
\[
F_{X_t, \dots, X_{t+j}} = F_{X_{t+h}, \dots, X_{t+h+j}}
\]
for all $t, j, h$.
\end{block}

\vspace{0.3cm}

\begin{itemize}
    \item Distribution invariant under time shifts
\end{itemize}
\end{frame}

%------------------------------------------------
\begin{frame}{Temporal Dependence}
\textbf{Goal: capture dependence across time}

\vspace{0.3cm}

\textbf{Problem:}
\begin{itemize}
    \item No meaningful arithmetic operations
    \item Classical ACF not applicable
\end{itemize}

\vspace{0.5cm}

\textbf{Solution:}
\begin{itemize}
    \item Work with joint probabilities
\end{itemize}

\vspace{0.3cm}

\begin{block}{Bivariate Probabilities}
\[
p_{ij}(h) = P(X_t = s_i, X_{t+h} = s_j), \quad  f_{ij}(h) = P(X_t \leq s_i, X_{t+h} \leq s_j)
\]
\end{block}

\begin{itemize}
    \item Independence $\Rightarrow p_{ij}(h) = p_i p_j \quad \text{and} \quad  f_{ij}(h) = f_i f_j $
\end{itemize}
\end{frame}

%------------------------------------------------
\begin{frame}{Cohen's $\kappa$}
\textbf{Idea: measure deviation from independence}

\vspace{0.5cm}

\[
\kappa(h)_{\mathrm{nom}} =
\frac{\sum_{j=0}^{m} (p_{jj}(h) - p_j^2)}
{1 - \sum_{i=0}^{m} p_i^2} \quad
\kappa(h)_{\mathrm{ord}} =
\frac{\sum_{j=0}^{m} (f_{jj}(h) - f_j^2)}
{1 - \sum_{i=0}^{m} f_i(1 - f_i)}
\]

\vspace{0.5cm}

\begin{itemize}
    \item $\kappa = 1$: perfect persistence
    \item $\kappa = 0$: independence
    \item $\kappa < 0$: tendency to switch states
\end{itemize}

\vspace{0.5cm}

\textbf{Interpretation:}
\begin{itemize}
    \item Probability of staying in same category
    \item Relative to independence benchmark
\end{itemize}
\end{frame}

%================================================
\section{Estimation — Complete Data}
%================================================

%------------------------------------------------
\begin{frame}{From Theory to Estimation}
\textbf{All quantities defined so far are population quantities.}

\vspace{0.4cm}

\begin{block}{Central question}
How do we estimate $\boldsymbol{f}$, IOV, and $\kappa_{\mathrm{ord}}(h)$ from data — and what are the statistical properties of these estimators?
\end{block}

\vspace{0.4cm}

\textbf{Two steps:}
\begin{enumerate}
    \item Estimate marginal CDF $\boldsymbol{f}$ via relative frequencies
    \item Derive distributions of plug-in estimators via the \textbf{delta method}
\end{enumerate}

\vspace{0.3cm}

$\Rightarrow$ Key tool: a CLT for dependent data under \textbf{mixing conditions}
\end{frame}

%------------------------------------------------
\begin{frame}{Estimation of Marginal Probabilities}
\textbf{Idea: use of indicator functions}

\vspace{0.3cm}

\[
Y_{t,i} = \mathbb{I}\{X_t = s_i\} \quad \text{and} \quad Z_{t,i} = \mathbb{I}\{X_t \leq s_i\}
\]
Then estimators are given by:
\begin{itemize}
    \item $\hat{p}_i = \frac{1}{n} \sum_{t=1}^n Y_{t,i}$
    \item $\hat{f}_i = \frac{1}{n} \sum_{t=1}^n Z_{t,i}$
\end{itemize}

\vspace{0.3cm}

Intuition:\\
$\Rightarrow$ we just count relative frequencies\\
From here we mainly look at ordinal data, but the nominal case is in many ways analogous.
\end{frame}

%------------------------------------------------
\begin{frame}{Mixing Assumptions}

\begin{block}{Idea}
The Central Limit Theorem (CLT) provides asymptotic normality and enables statistical inference.
\end{block}

\begin{alertblock}{Problem}
Classical CLT results assume independence, but in practice data may exhibit serial dependence.
\end{alertblock}

\begin{block}{Solution}
To handle dependence, we impose \textbf{mixing conditions}:
\begin{itemize}
    \item $\alpha$-mixing quantifies the decay of dependence between past and future observations
    \item Exponential decay of the mixing coefficients allows a CLT for dependent data
\end{itemize}
\end{block}
\end{frame}

%------------------------------------------------
\begin{frame}{CLT under Serial Dependence}
    \begin{block}{Result}
    \[\sqrt{n}(\hat{f} - f) \Rightarrow \mathcal{N}(0, \boldsymbol{\Sigma})\]
    where the elements of $\boldsymbol{\Sigma}$ are given by:
    \[\sigma_{ij} = \underbrace{\left(f_{\min\{i,j\}}- f_i f_j\right)}_{\text{lag}(0)} + \sum_{h=1}^{\infty} \underbrace{\left(f_{ij}(h) + f_{ji}(h) - 2f_if_j\right)}_{\text{lag}(h)}.\]
    Note that $\text{lag}(h) \rightarrow 0$ as $h \rightarrow \infty$.
    \end{block}
\end{frame}

%------------------------------------------------
\begin{frame}{CLTs for Marginal Estimators}
    \begin{block}{Idea}
        IOV is a function of the marginal probabilities \\
        $\Rightarrow$ apply the \textbf{delta method} to obtain asymptotic distributions 
    \end{block}
    \begin{block}{Results}
        \begin{itemize}
            \item Asymptotic variance: 
            \[\sqrt{n}\,\mathrm{Var}(\widehat{\mathrm{IOV}}) \overset{\mathrm{asy}}{=} \frac{16}{m^2}\sum_{i,j = 0}^{m-1} (1 - 2f_i)(1 - 2f_j) \sigma_{ij}\]
            \item Bias to order $n^{-1}$:
            \[n\, \mathrm{Bias}(\widehat{\mathrm{IOV}}) = -\frac{4}{m} \sum_{i = 0}^{m-1}  \sigma_{ii}\]
        \end{itemize}
    \end{block}
\end{frame}

%================================================
\section{Missing Data}
%================================================

%------------------------------------------------
\begin{frame}{Why Does Missingness Matter?}

\textbf{Missing observations are ubiquitous in practice:}
\begin{itemize}
    \item Panel surveys: participants drop out over time
    \item Administrative data: recording gaps or system failures
    \item Sensor data: transmission errors, device downtime
\end{itemize}

\vspace{0.4cm}

\textbf{Key problem: naive estimators are biased}

\begin{itemize}
    \item Simply ignoring missing observations distorts relative frequencies
    \item $\Rightarrow$ biased estimates of $\boldsymbol{f}$, IOV, and $\kappa_{\mathrm{ord}}(h)$
\end{itemize}

\vspace{0.4cm}

\textbf{Goal:}
\begin{itemize}
    \item Correct for missingness — under suitable assumptions
    \item Derive valid asymptotic theory for the corrected estimators
\end{itemize}
\end{frame}

%------------------------------------------------
\begin{frame}{Amplitude Modulation}
    \begin{block}{Idea}
        Missingness can be modeled as a binary process $(O_t)$, that \textit{modulates} the \textit{amplitude} of the actual process: 
        \[O_t = 1 \text{ if } X_t \text{ is observed and } O_t = 0 \text{ if } X_t \text{ is missing}\] 
        The observed process is then given by $(O_t X_t)$.
    \end{block}
    \begin{block}{Assumptions}
        \begin{itemize}
            \item $(O_t)$ is stationary with $P(O_t = 1) = \pi > 0$ and joint moments $E[O_t O_{t+h}] = \pi(h)$
            \item $(O_t)$ is $\alpha$-mixing with exponential decay
            \item $(O_t)$ and $(X_t)$ are \textbf{independent}
        \end{itemize}
    \end{block}    
\end{frame}

%------------------------------------------------
\begin{frame}{Missingness Mechanisms}
    Originally developed for cross-sectional data, but also relevant for time series:
    \begin{itemize}  
        \item \textbf{MCAR:} Missing Completely At Random
        \[P(O_t = 1| X_t, X_{t-1}, \ldots)\ = P(O_t = 1)\]
        \item \textbf{MAR:} Missing At Random
        \[P(O_t = 1| X_t, X_{t-1}, \ldots)\ = P(O_t = 1| X_{t-1}, \ldots)\]
        \item \textbf{MNAR:} Missing Not at Random
        \[P(O_t = 1| X_t, X_{t-1}, \ldots)\ \neq P(O_t = 1| X_{t-1}, \ldots)\]
    \end{itemize}
    From here on, unless stated otherwise, we assume MCAR.
\end{frame}

%------------------------------------------------
\begin{frame}{Naive vs.\ Corrected Estimators}
    \begin{block}{Naive estimator}
        \[\hat{\boldsymbol{\mathfrak{f}}} = \frac{1}{n} \sum_{t=1}^n O_t \boldsymbol{Z}_t\]
        $\Rightarrow$ biased and inconsistent: $E[\hat{\mathfrak{f}}_i] = \pi f_i \neq f_i$
    \end{block}
    \begin{block}{Corrected estimator}
        \[\hat{\boldsymbol{f}}^* = \frac{\sum_{t=1}^n O_t \boldsymbol{Z}_t}{\sum_{t=1}^n O_t}\]
        $\Rightarrow$ divide by number of observed time points, not by $n$
    \end{block}
\end{frame}

%------------------------------------------------
\begin{frame}{Asymptotic Properties of the Corrected Estimator}
    \begin{block}{Result}
        \[\sqrt{n}(\hat{\boldsymbol{f}}^* - \boldsymbol{f}) \Rightarrow \mathcal{N}(0, \boldsymbol{\Sigma}^*)\]
        where the elements of $\boldsymbol{\Sigma}^*$ are given by:
        \[\sigma^*_{ij} = \frac{1}{\pi} (f_{\mathrm{min}\{i,j\}}- f_i f_j) + \frac{1}{\pi^2} \sum_{h = 1}^{\infty} \pi(h)(f_{ij}(h) + f_{ji}(h) - 2f_if_j)\]
    \end{block}

    \vspace{0.3cm}
    \textbf{Remarks:}
    \begin{itemize}
        \item Inflation factor $1/\pi$ at lag 0: fewer effective observations
        \item Serial dependence in $(O_t)$ enters via $\pi(h) = E[O_t O_{t+h}]$
        \item Reduces to $\boldsymbol{\Sigma}$ when $\pi = 1$ (no missing data)
    \end{itemize}
\end{frame}

%------------------------------------------------
\begin{frame}{Marginal Statistics under Missingness}
    \begin{block}{Strategy}
        Replace $\boldsymbol{\Sigma}$ by $\boldsymbol{\Sigma}^*$ in all delta-method results
    \end{block}

    \vspace{0.3cm}

    \textbf{IOV} (nonlinear in $\boldsymbol{f}$, delta method applies):
    \[
    \sqrt{n}\,\mathrm{Var}\bigl(\mathrm{IOV}(\hat{\boldsymbol{f}}^*)\bigr) \overset{\mathrm{asy.}}{=}
    \frac{16}{m^2} \sum_{i,j=0}^{m-1}(1-2f_i)(1-2f_j)\,\sigma^*_{i,j}
    \]
    \[
    \mathrm{Bias}\bigl(\mathrm{IOV}(\hat{\boldsymbol{f}}^*)\bigr) = -\frac{4}{mn}\sum_{i=0}^{m-1}\sigma^*_{i,i}
    \]

    \vspace{0.3cm}

    \textbf{Skew} (linear in $\boldsymbol{f}$, Hessian vanishes):
    \begin{itemize}
        \item Estimator is \textbf{unbiased}
        \item Asymptotic variance analogous, replacing $\boldsymbol{\Sigma}$ by $\boldsymbol{\Sigma}^*$
    \end{itemize}
\end{frame}

%================================================
\section{Serial Dependence under Missingness}
%================================================

%------------------------------------------------
\begin{frame}{Estimating Bivariate Probabilities}
\textbf{Need to estimate $f_{ij}(h) = P(X_t \leq s_i, X_{t+h} \leq s_j)$}

\vspace{0.3cm}

\textbf{Naive estimator:}
\[
\hat{\mathfrak{f}}_{ij}(h) = \frac{1}{n-h}\sum_{t=h+1}^{n} O_t\, O_{t-h}\, Z_{t,i}\, Z_{t-h,j}
\]

\textbf{Problem:} $E[\hat{\mathfrak{f}}_{ij}(h)] = \pi(h) \cdot f_{ij}(h) \neq f_{ij}(h)$

\vspace{0.4cm}

\textbf{Corrected estimator:}
\[
\hat{f}^*_{ij}(h) = \frac{\frac{1}{n-h}\sum_{t=h+1}^{n} O_t\, O_{t-h}\, Z_{t,i}\, Z_{t-h,j}}{\frac{1}{n-h}\sum_{t=h+1}^{n} O_t\, O_{t-h}} =: \hat{\mathfrak{f}}_{ij}(h) \Big/ \overline{OO}_h
\]

$\Rightarrow$ divide by the proportion of \textbf{jointly observed} pairs at lag $h$
\end{frame}

%------------------------------------------------
\begin{frame}{Estimating Cohen's $\kappa$ under Missingness}
\textbf{Plug-in estimator:}
\[
\hat{\kappa}_{\mathrm{ord}}(h) = \frac{\sum_{j=0}^{m-1}\left(\hat{f}_{jj}^*(h) - (\hat{f}^*_j)^2\right)}{\sum_{i=0}^{m-1}\hat{f}^*_i(1-\hat{f}_i^*)}
\]

\vspace{0.4cm}

\textbf{Joint asymptotic distribution via delta method:}
\begin{itemize}
    \item Define augmented vector $\hat{\boldsymbol{f}}^{(h)*} = (\hat{f}^*_0, \ldots, \hat{f}^*_{m-1}, \hat{f}^*_{00}(h), \ldots, \hat{f}^*_{m-1,m-1}(h))^\top$
    \item CLT holds for the sample mean of the underlying binary process $\boldsymbol{Z}^{(h)*}_t$
    \item Apply delta method to $\kappa_{\mathrm{ord}}$
\end{itemize}

\vspace{0.3cm}
\textbf{Key challenge:} long-run covariance involves third- and fourth-order mixed moments of $(X_t)$ — not closed-form in general.\\
$\Rightarrow$ focus on the \textbf{null distribution} first
\end{frame}

%------------------------------------------------
\begin{frame}{Asymptotic Null Distribution of $\hat{\kappa}_{\mathrm{ord}}(h)$}
\textbf{Under $H_0$: $(X_t)$ is i.i.d.}, $f_{ij}(h) = f_i f_j$, and higher-order moments factor.

\vspace{0.3cm}

\begin{block}{Theorem}
Under $H_0$ and MCAR,
\[
\sqrt{n}\,\hat{\kappa}_{\mathrm{ord}}(h) \xrightarrow{d} \mathcal{N}(0, \sigma^2_\kappa),
\]
where
\[
\sigma^2_\kappa =
\frac{\dfrac{1}{\pi(h)}\displaystyle\sum_{i,j=0}^{m-1}
(f_{\min\{i,j\}}-f_if_j)
\!\left(f_{\min\{i,j\}} + \left(1+\dfrac{2\pi(h,2h)}{\pi(h)} - \dfrac{4\pi(h)}{\pi}\right)f_if_j\right)}
{\left(\displaystyle\sum_{k=0}^{m-1}f_k(1-f_k)\right)^2}
\]
with $\pi(h,2h) = E[O_t O_{t+h} O_{t-h}]$.
\end{block}
\end{frame}

%------------------------------------------------
\begin{frame}{Bias Correction for $\hat{\kappa}_{\mathrm{ord}}(h)$}

\textbf{Second-order Taylor expansion yields the finite-sample bias under $H_0$:}
\[
\mathrm{Bias}\!\left(n\hat{\kappa}_{\mathrm{ord}}(h)\right) \approx -\frac{1}{\pi}
\]

\vspace{0.4cm}

\textbf{Bias-corrected estimator:}
\[
\hat{\kappa}^{\mathrm{bc}}_{\mathrm{ord}}(h) := \hat{\kappa}_{\mathrm{ord}}(h) + \frac{1}{n\pi}
\]

\vspace{0.4cm}

\textbf{Remarks:}
\begin{itemize}
    \item Bias increases as $\pi \downarrow$ (more missingness)
    \item Under complete data ($\pi = 1$): reduces to the known $-1/n$ bias
    \item Closed-form correction — no resampling needed
\end{itemize}
\end{frame}

%------------------------------------------------
\begin{frame}{Consistency under $H_A$}

\textbf{Under $H_A$:} $(X_t)$ is not i.i.d., so $\kappa_{\mathrm{ord}}(h) \neq 0$ for some $h$.

\vspace{0.4cm}

\begin{block}{Result}
$\hat{\kappa}_{\mathrm{ord}}(h)$ is a continuous function of $\hat{f}^*_i$ and $\hat{f}^*_{ii}(h)$.\\
By the \textbf{continuous mapping theorem}:
\[
\hat{\kappa}_{\mathrm{ord}}(h) \xrightarrow{p} \kappa_{\mathrm{ord}}(h) \neq 0,
\]
and therefore
\[
\sqrt{n}\,\hat{\kappa}_{\mathrm{ord}}(h) \xrightarrow{p} \infty.
\]
\end{block}

\vspace{0.3cm}
$\Rightarrow$ The test is \textbf{consistent}: rejection probability $\to 1$ under $H_A$.
\end{frame}

%================================================
\section{Simulation Study}
%================================================

%------------------------------------------------
\begin{frame}{Simulation Goals and Setup}
\textbf{Three goals:}
\begin{enumerate}
    \item Assess finite-sample accuracy of asymptotic bias formulas
    \item Examine convergence to asymptotic normality
    \item Study robustness to violations of the MCAR assumption
\end{enumerate}

\vspace{0.4cm}

\textbf{Data generating process:} BinAR(1) model
\begin{itemize}
    \item Flexible ordinal time series with controllable serial dependence
    \item Stationary marginal distribution: $X_t \sim \mathrm{Binomial}(m, p)$
    \item Parameter $r \in (-1,1)$ controls autocorrelation ($r=0$: i.i.d.)
    \item All quantities (IOV, $\kappa_{\mathrm{ord}}$) available in \textbf{closed form}
\end{itemize}
\end{frame}

%------------------------------------------------
\begin{frame}{Simulation Scenarios}
\textbf{Two scenarios chosen to differ qualitatively:}

\vspace{0.3cm}

\begin{columns}
\begin{column}{0.48\textwidth}
\textbf{Scenario A}\\
$m=3$, $p=0.20$, $r=0.35$
\begin{itemize}
    \item Few categories
    \item Skewed marginal distribution
    \item Moderate serial dependence
\end{itemize}
\end{column}
\begin{column}{0.48\textwidth}
\textbf{Scenario B}\\
$m=10$, $p=0.45$, $r=0.50$
\begin{itemize}
    \item Many categories
    \item Near-symmetric distribution
    \item Strong serial dependence
\end{itemize}
\end{column}
\end{columns}

\vspace{0.4cm}

\textbf{Design:}
\begin{itemize}
    \item Observation rates: $\pi \in \{1, 0.75\}$
    \item Sample sizes: $n \in \{50, 100, 250, 500, 1000\}$
    \item $B = 1{,}000$ Monte Carlo replications
\end{itemize}
\end{frame}

%------------------------------------------------
\begin{frame}{Missingness Mechanisms in the Simulation}
\textbf{Four observation processes $(O_t)$ are studied:}

\begin{enumerate}
    \item \textbf{Independent MCAR}: $O_t \overset{\mathrm{i.i.d.}}{\sim} \mathrm{Bernoulli}(\pi)$ \\ {\small Baseline; satisfies all theoretical assumptions.}
    
    \vspace{0.3cm}
    
    \item \textbf{Dependent MCAR}: $(O_t)$ follows BinAR(1) with $m=1$ \\ {\small Serial dependence in missingness; MCAR still holds marginally.}
    
    \vspace{0.3cm}
    
    \item \textbf{MNAR — increasing}: $\pi_t = \pi_{\mathrm{low}} + (\pi_{\mathrm{high}} - \pi_{\mathrm{low}})\frac{X_t}{m}$ \\ {\small Higher categories more likely to be observed.}
    
    \vspace{0.3cm}
    
    \item \textbf{MNAR — decreasing}: $\pi_t = \pi_{\mathrm{high}} - (\pi_{\mathrm{high}} - \pi_{\mathrm{low}})\frac{X_t}{m}$ \\ {\small Lower categories more likely to be observed.}
\end{enumerate}
\end{frame}

%------------------------------------------------
\begin{frame}{Performance Measures}
\textbf{For each estimator $\hat{\theta}$, we compute across $B$ replications:}

\vspace{0.3cm}

\begin{itemize}
    \item \textbf{Mean bias:} $\frac{1}{B}\sum_{b=1}^B (\hat{\theta}_b - \theta)$ \\
    {\small Compared against theoretical $O(n^{-1})$ formula.}
    
    \vspace{0.3cm}
    
    \item \textbf{$n$-scaled bias:} $n \cdot \frac{1}{B}\sum_{b=1}^B (\hat{\theta}_b - \theta)$ \\
    {\small Should converge to the theoretical bias limit as $n \to \infty$.}
    
    \vspace{0.3cm}
    
    \item \textbf{Standard deviation:} compared against asymptotic SD $\sqrt{\sigma^2_\theta/n}$
\end{itemize}

\vspace{0.4cm}

\textbf{Theoretical benchmarks:}
\begin{itemize}
    \item IOV bias limit: $-\frac{4}{mn}\,\mathrm{tr}(\boldsymbol{\Sigma}^*)$
    \item $\kappa$ bias limit (under $H_0$): $-\frac{1}{n\pi}$
\end{itemize}
\end{frame}

%================================================
\section{Summary}
%================================================

%------------------------------------------------
\begin{frame}{Summary}
\begin{itemize}
    \item Categorical time series require new tools:
    \begin{itemize}
        \item Nominal vs.\ ordinal $\Rightarrow$ PMF vs.\ CDF, mode vs.\ median
    \end{itemize}
    
    \vspace{0.2cm}
    
    \item Marginal estimation:
    \begin{itemize}
        \item Relative frequency estimators + CLT under $\alpha$-mixing
        \item Delta method for IOV and skew
    \end{itemize}
    
    \vspace{0.2cm}

    \item Missingness:
    \begin{itemize}
        \item Amplitude modulation framework under MCAR
        \item Corrected estimators $\hat{f}^*$ and $\hat{f}^*_{ij}(h)$
        \item Full asymptotic theory carries over via delta method
    \end{itemize}
    
    \vspace{0.2cm}

    \item Serial dependence under missingness:
    \begin{itemize}
        \item Asymptotic null distribution of $\hat{\kappa}_{\mathrm{ord}}(h)$
        \item Closed-form bias correction $\hat{\kappa}^{\mathrm{bc}}_{\mathrm{ord}}(h) = \hat{\kappa}_{\mathrm{ord}}(h) + \frac{1}{n\pi}$
        \item Consistency under $H_A$
    \end{itemize}
    
    \vspace{0.2cm}

    \item Simulation study validates finite-sample behavior and robustness
\end{itemize}
\end{frame}

\end{document}
